\subsection{Supplementary GA heuristic baseline}
\label{sec:ga-baseline}
Sarhan et al. study a genetic algorithm for radio labeling
\cite{sarhan2026genetic}; radio-labeling constraints are not interchangeable
with the two-distance $L(h,k)$ constraints considered here. We adapt the
permutation-search idea only. A chromosome orders the vertices; the decoder
assigns each vertex the smallest nonnegative label satisfying threshold $h$
against already labeled neighbors and threshold $k$ against already labeled
vertices at distance two. Pairs at larger distance, including across connected
components, have no constraint. Each decoded chromosome is feasible, and its
fitness is the labeling span.

The fixed configuration uses a population of 40, OX1 crossover with
probability 0.9, segment inversion with probability 0.05, tournament
selection of size three, and elitism. Ten independent algorithmic seeds are
run on each of the same 117 graph/parameter tasks used in the progress-v2
SAT--ILP study; each GA run decodes 1,000 permutations. As a control, a
random-order multi-start method uses the same decoder, evaluation budget and
paired seeds. A graph-insertion-order greedy labeling is also retained as a
single-pass baseline. Settings are fixed rather than tuned on a held-out set.

The optimum for every task is independently replayed from the proof-closure
package, and each heuristic witness is rechecked by the separate validator.
At least one of ten GA runs attains the certified optimum on 95/117 tasks,
compared with 85/117 for random-order search. Aggregating the median of
seed-paired span differences per task gives 21 GA wins, 96 ties and no losses;
15 of the 18 ER $L(3,2)$ tasks are GA wins. This descriptive difference is
concentrated in that family. It does not establish general superiority of
genetic search. The signed span difference $\lambda_{\rm SAT}-\lambda_{\rm
GA}$ is reported only when SAT was optimal in all three progress-v2 repeats;
the standard nonnegative GA optimality gap is instead
$\lambda_{\rm GA}-\lambda^*$. GA and random-order runtimes exclude shared
constraint preprocessing and witness validation, so they are not compared
directly with SAT/ILP wall times.

\begin{table}[htbp]
\centering\small
\caption{Fixed-budget GA and random-order outcomes by family and parameter.}
\label{tab:ga-baseline}
\begin{tabular}{llrrrrrr}
\toprule
Family & $(h,k)$ & Cases & GA opt. & Rand. opt. & GA wins & Ties & GA losses\\
\midrule
BA & $(1,1)$ & 9 & 9 & 9 & 0 & 9 & 0\\
BA & $(2,1)$ & 9 & 9 & 9 & 0 & 9 & 0\\
BA & $(3,2)$ & 9 & 9 & 9 & 2 & 7 & 0\\
ER & $(1,1)$ & 18 & 15 & 14 & 1 & 17 & 0\\
ER & $(2,1)$ & 18 & 11 & 12 & 2 & 16 & 0\\
ER & $(3,2)$ & 18 & 12 & 3 & 15 & 3 & 0\\
PxP & $(1,1)$ & 3 & 1 & 1 & 0 & 3 & 0\\
PxP & $(2,1)$ & 3 & 1 & 1 & 1 & 2 & 0\\
PxP & $(3,2)$ & 3 & 1 & 1 & 0 & 3 & 0\\
tree & $(1,1)$ & 9 & 9 & 9 & 0 & 9 & 0\\
tree & $(2,1)$ & 9 & 9 & 8 & 0 & 9 & 0\\
tree & $(3,2)$ & 9 & 9 & 9 & 0 & 9 & 0\\
\bottomrule
\end{tabular}

\tablenote{GA opt. and Random opt. count graph/parameter tasks on which the
best of ten seeds attains the certified optimum. Win/tie/loss compares the
median of seed-paired GA-minus-random span differences within each task;
negative favors GA. Repeated seeds do not add graph instances.}
\end{table}
